\subsection{Actual weighted norm forms over every complete base field}
\label{subsec:general-weighted-norms}

The determinant polynomial contains more information than its values on the
original field. This section uses that polynomial information to prove the
weighted norm identity also when the completed residue field is finite.

\begin{lemma}[Whole norm polynomials and base change]
\label{spectrum:whole-norm-form-base-change}
\lean{EnvelopeProperty.map_normForm_baseChange,
EnvelopeProperty.normForm_algEquiv,
EnvelopeProperty.normForm_eq_prod_of_algEquiv}
\leanok
\uses{spectrum:algebra-norm-base-change}
Let $B$ be finite free over a commutative ring $A$ and let $K$ be any
commutative $A$-algebra. The actual determinant norm polynomial satisfies
\[
 \operatorname{map}_{A\to K}N_{B/A}\Bigl(\sum_j X_jb_j\Bigr)
 =N_{K\otimes_A B/K}\Bigl(\sum_j X_j(1\otimes b_j)\Bigr).
\]
Algebra isomorphisms preserve this whole polynomial. A genuine product
isomorphism to finite free factors factors it into the product of their
whole determinant norm polynomials.
\end{lemma}
\begin{proof}
Choose a finite free basis of $B$. In the scalar-extended basis and then its
polynomial extension, the multiplication matrix of the linear form is
obtained by applying the coefficient map to every entry of the original
polynomial multiplication matrix. Determinants commute with this ring map,
proving coefficient base change as an equality of polynomials.
An algebra isomorphism induces the actual polynomial tensor isomorphism;
invariance of determinant norm under that isomorphism proves the second
assertion. For a product, the tensor product with the polynomial ring
commutes with the finite product. The norm of this product algebra is the
product of the factor norms, yielding the whole-polynomial factorization.
No polynomial interpolation is used in these three identities.
\end{proof}

\begin{theorem}[Weighted field norm form without infinitude]
\label{spectrum:weighted-field-norm-without-infinitude}
\lean{EnvelopeProperty.wgaussValuation_normForm_field}
\leanok
\uses{spectrum:whole-norm-form-base-change,spectrum:complete-norm-extension}
Let $H$ be a complete ultrametric normed field, $L/H$ a finite separable
extension, and $w$ its extending real valuation. For finitely many $b_j\in L$
and arbitrary $r_j\ge0$,
\[
 \left\|N_{L/H}\Bigl(\sum_j X_jb_j\Bigr)\right\|_r
 =\left(\max_j w(b_j)r_j\right)^{[L:H]}.
\]
There is no infinitude hypothesis on $H$.
\end{theorem}
\begin{proof}
Let $\Omega$ be an algebraic closure of $H$ and choose its extending
valuation $w_\Omega$. The actual finite separable base change
$\Omega\otimes_H L$ is formally etale and is a product of copies of
$\Omega$ indexed by its prime spectrum. Each factor projection, restricted
along $L\to\Omega\otimes_H L$, is an $H$-embedding $\tau_m:L\to\Omega$.
Complete-field uniqueness gives $w_\Omega(\tau_m(b))=w(b)$ for every $b$.

Evaluate the scalar-extended actual norm polynomial at any
$t\in\Omega^{\{j\}}$. The determinant base-change and evaluation identity
makes its value the norm of $\sum_j t_j(1\otimes b_j)$ in the actual split
algebra. The split algebra computes this norm as
$\prod_m\sum_j t_j\tau_m(b_j)$. Since the algebraically closed field
$\Omega$ is infinite, these evaluations prove equality of the whole
polynomial with the product of the indicated linear polynomials.
Infinitude of the original field is never needed.

Each linear factor has weighted Gauss norm $\max_j w(b_j)r_j$.
Multiplicativity gives one such factor for each prime of the split algebra. Its
number of factors equals $[L:H]$: dimension is preserved by tensor base
change and the product of that many copies of $\Omega$ has that dimension.
Coefficient valuation compatibility returns the weighted norm to $H$.
No complete or normed field structure on $\Omega$ is used; only its
extending valuation is needed.
\end{proof}

\begin{theorem}[Weighted norm forms for products of separable fields]
\label{spectrum:weighted-product-norm-without-infinitude}
\lean{EnvelopeProperty.wgaussValuation_normForm_of_algEquiv}
\leanok
\uses{spectrum:whole-norm-form-base-change,spectrum:weighted-field-norm-without-infinitude}
Suppose a commutative $H$-algebra has a genuine isomorphism
$E\simeq\prod_iL_i$ with finite separable field extensions, each carrying
its extending valuation $w_i$. Then
\[
 \left\|N_{E/H}\Bigl(\sum_j X_jb_j\Bigr)\right\|_r
 =\prod_i\left(\max_jw_i(b_{ij})r_j\right)^{[L_i:H]}.
\]
The complete field $H$ may be finite or trivially normed.
\end{theorem}
\begin{proof}
The actual product isomorphism factors the entire determinant norm
polynomial by the first lemma. Weighted Gauss multiplicativity reduces
its value to the product of the factor values. Apply the preceding field
formula to every finite separable factor. This proves the polynomial
identity without replacing the norm polynomial by an unspecified
polynomial with the same values on $H$.
\end{proof}

\begin{theorem}[Weighted finite etale norm identity at every seminorm point]
\label{spectrum:weighted-etale-norm-without-infinitude}
\lean{EnvelopeProperty.Berkovich.Spectrum.wgaussValuation_map_normForm_of_etale,
EnvelopeProperty.Berkovich.Spectrum.wgaussValuation_normForm_of_etale}
\leanok
\uses{spectrum:whole-norm-form-base-change,spectrum:weighted-product-norm-without-infinitude,
spectrum:etale-fibre-multiplicities}
Let $A\to B$ be finite free etale, $x\in\mathcal S(A)$, and let
$N_{B/A}(\sum_j X_jb_j)=\sum_\alpha c_\alpha X^\alpha$. For all $r_j\ge0$,
\[
 \max_\alpha |c_\alpha|_x r^\alpha
 =\prod_{y\text{ over }x}\left(\max_j|b_j|_y r_j\right)^{m_y},
 \qquad m_y=[(H(x)\otimes_A B)/\mathfrak m_y:H(x)].
\]
This includes finite completed residue fields and arbitrary prescribed
coefficient seminorms.
\end{theorem}
\begin{proof}
Whole-polynomial coefficient base change identifies the image of the
original determinant norm polynomial in $H(x)[X]$ with the actual norm
form of $E_x=H(x)\otimes_A B$. The reduced Artinian etale decomposition
of $E_x$ gives genuine finite separable maximal quotient fields.
The actual completed tensor maps to $H(y)$ descend to these quotients,
provide valuations extending the norm of $H(x)$, and take value $|b_j|_y$
on the component of $1\otimes b_j$.

Apply the preceding product-field formula to this actual norm form.
The genuine maximal-kernel classification of completed tensor fibres
reindexes every factor by the corresponding point $y$ and identifies
its exponent with the defined residue-degree multiplicity. Preservation
of the coefficient seminorm under $A\to H(x)$ gives the stated original
coefficient formula. These steps use equality on whole polynomials and
whole tensor points. Etaleness is load-bearing: a nonetale norm formula
would also require local-length factors, which this theorem does not
assert.
\end{proof}

The older evaluation-characterized polynomial lemmas still require an
infinite original field: over a finite field, different polynomials can
agree at every field point. The new theorem applies to the actual
determinant norm polynomial and removes that restriction by using its
proved coefficient base-change identity.
